\chapter{Mathematical Derivations}
\label{chap:appendix-derivations}

This appendix collects derivations that are useful for understanding the topology-related quantities written by FEMaster. The main reference chapters describe how to use the DSL commands. This appendix gives the mathematical background for compliance, density sensitivities, and orientation sensitivities.

\section{Element stiffness}

For a linear elastic finite element, the unmodified element stiffness is
\[
K_{e,0} = \int_{\Omega_e} B^T D B \, d\Omega .
\]
\(B\) is the strain-displacement matrix, \(D\) is the constitutive matrix, and \(\Omega_e\) is the physical element domain. With an isoparametric mapping from natural coordinates \(r=(\xi,\eta,\zeta)\), the integral is evaluated as
\[
K_{e,0}
=
\int_{\Omega'_e}
B(r)^T D B(r)\,
\lvert \det J(r)\rvert \, d\xi\,d\eta\,d\zeta .
\]
For numerical integration with quadrature points \(r_q\) and weights \(w_q\), this becomes
\[
K_{e,0}
\approx
\sum_{q=1}^{Q}
w_q B(r_q)^T D B(r_q)\,
\lvert \det J(r_q)\rvert .
\]

\section{Compliance}

For a linear static analysis,
\[
K u = f .
\]
The external work measure commonly used for topology optimization is the compliance
\[
C = f^T u .
\]
Using the equilibrium equation \(f=Ku\), this can also be written as
\[
C = u^T K u .
\]
Because the global stiffness matrix is assembled from element stiffness matrices, compliance can be decomposed into element contributions:
\[
C = \sum_{e=1}^{n_e} u_e^T K_e u_e .
\]
This expression is the basis for element-wise topology sensitivities. \(u_e\) is the element displacement vector extracted from the global solution.

\section{Density penalization}

In a density-based topology formulation, an element density variable \(\rho_e\) scales the stiffness contribution. With a penalization exponent \(p\),
\[
K_e(\rho_e) = \rho_e^p K_{e,0}.
\]
The exponent \(p\) makes intermediate densities less efficient and therefore encourages designs closer to a solid-void distribution.

The element compliance contribution is
\[
C_e(\rho_e) = u_e^T K_e(\rho_e) u_e
=
\rho_e^p u_e^T K_{e,0} u_e .
\]
If \(u_e\) is held fixed for the local sensitivity expression, the derivative is
\[
\frac{\partial C_e}{\partial \rho_e}
=
p\rho_e^{p-1} u_e^T K_{e,0} u_e .
\]
Depending on the optimization convention, the objective derivative may be reported with the opposite sign because minimizing compliance through the equilibrium equation often uses the adjoint result
\[
\frac{dC}{d\rho_e}
=
-u_e^T \frac{\partial K_e}{\partial \rho_e}u_e
=
-p\rho_e^{p-1}u_e^T K_{e,0}u_e .
\]
The sign convention should therefore be interpreted together with the optimizer interface and the field name written by the analysis.

\section{Orientation-dependent stiffness}

For orientation-based topology or material steering, the constitutive matrix is rotated before the element stiffness is evaluated. Let the material orientation be represented by three angles
\[
\alpha = (\alpha_1,\alpha_2,\alpha_3).
\]
The elementary rotation matrices are
\[
R_1 =
\begin{bmatrix}
1 & 0 & 0\\
0 & \cos\alpha_1 & -\sin\alpha_1\\
0 & \sin\alpha_1 & \cos\alpha_1
\end{bmatrix},
\]
\[
R_2 =
\begin{bmatrix}
\cos\alpha_2 & 0 & \sin\alpha_2\\
0 & 1 & 0\\
-\sin\alpha_2 & 0 & \cos\alpha_2
\end{bmatrix},
\qquad
R_3 =
\begin{bmatrix}
\cos\alpha_3 & -\sin\alpha_3 & 0\\
\sin\alpha_3 & \cos\alpha_3 & 0\\
0 & 0 & 1
\end{bmatrix}.
\]
For the convention used here, the combined rotation is
\[
R = R_1R_2R_3 .
\]
The multiplication order must be used consistently. Different conventions can describe the same physical orientation with different angle triples.

\section{Voigt transformation}

The strain vector uses engineering shear strains,
\[
\varepsilon =
\begin{bmatrix}
\varepsilon_{xx} &
\varepsilon_{yy} &
\varepsilon_{zz} &
2\varepsilon_{yz} &
2\varepsilon_{xz} &
2\varepsilon_{xy}
\end{bmatrix}^T ,
\]
and the stress vector is
\[
\sigma =
\begin{bmatrix}
\sigma_{xx} &
\sigma_{yy} &
\sigma_{zz} &
\sigma_{yz} &
\sigma_{xz} &
\sigma_{xy}
\end{bmatrix}^T .
\]
The strain transformation is written as
\[
\varepsilon_m = T(R)\varepsilon_g ,
\]
where \(T\) is the \(6\times6\) Voigt transformation matrix built from the entries of \(R\). The rotated constitutive matrix used in global strain coordinates is
\[
D_\mathrm{rot}=T^TDT .
\]
The orientation-dependent element stiffness is therefore
\[
K_e(\rho_e,\alpha_e)
=
\rho_e^p
\int_{\Omega_e}
B^T D_\mathrm{rot}(\alpha_e) B\,d\Omega .
\]
With quadrature,
\[
K_e(\rho_e,\alpha_e)
\approx
\rho_e^p
\sum_{q=1}^{Q}
w_q B(r_q)^T T(\alpha_e)^T D T(\alpha_e)B(r_q)
\lvert \det J(r_q)\rvert .
\]

\section{Orientation sensitivity}

For one element, the compliance contribution is
\[
C_e
=
\rho_e^p
\sum_{q=1}^{Q}
w_q
\varepsilon_q^T
T^TDT
\varepsilon_q
\lvert \det J_q\rvert ,
\qquad
\varepsilon_q = B(r_q)u_e .
\]
Only \(T\) depends on the orientation angles, so the derivative with respect to angle \(\alpha_j\) is
\[
\frac{\partial C_e}{\partial \alpha_j}
=
\rho_e^p
\sum_{q=1}^{Q}
w_q
\varepsilon_q^T
\left(
\frac{\partial T^T}{\partial \alpha_j}DT
+
T^TD\frac{\partial T}{\partial \alpha_j}
\right)
\varepsilon_q
\lvert \det J_q\rvert .
\]
If \(D\) is symmetric, the scalar expression can also be evaluated as
\[
\frac{\partial C_e}{\partial \alpha_j}
=
2\rho_e^p
\sum_{q=1}^{Q}
w_q
\varepsilon_q^T
T^TD\frac{\partial T}{\partial \alpha_j}
\varepsilon_q
\lvert \det J_q\rvert ,
\]
provided the same transformation convention is used.

\section{Rotation derivatives}

The derivatives of the elementary rotation matrices are
\[
\frac{\partial R_1}{\partial \alpha_1}
=
\begin{bmatrix}
0 & 0 & 0\\
0 & -\sin\alpha_1 & -\cos\alpha_1\\
0 & \cos\alpha_1 & -\sin\alpha_1
\end{bmatrix},
\]
\[
\frac{\partial R_2}{\partial \alpha_2}
=
\begin{bmatrix}
-\sin\alpha_2 & 0 & \cos\alpha_2\\
0 & 0 & 0\\
-\cos\alpha_2 & 0 & -\sin\alpha_2
\end{bmatrix},
\]
\[
\frac{\partial R_3}{\partial \alpha_3}
=
\begin{bmatrix}
-\sin\alpha_3 & -\cos\alpha_3 & 0\\
\cos\alpha_3 & -\sin\alpha_3 & 0\\
0 & 0 & 0
\end{bmatrix}.
\]
For \(R=R_1R_2R_3\), the total rotation derivatives are
\[
\frac{\partial R}{\partial \alpha_1}
=
\frac{\partial R_1}{\partial \alpha_1}R_2R_3,
\]
\[
\frac{\partial R}{\partial \alpha_2}
=
R_1\frac{\partial R_2}{\partial \alpha_2}R_3,
\]
\[
\frac{\partial R}{\partial \alpha_3}
=
R_1R_2\frac{\partial R_3}{\partial \alpha_3}.
\]
The derivative \(\partial T/\partial\alpha_j\) is obtained by differentiating every entry of \(T(R)\) with the product rule. For example, an entry \(R_{a b}R_{c d}\) contributes
\[
\frac{\partial}{\partial\alpha_j}
\left(R_{a b}R_{c d}\right)
=
\frac{\partial R_{a b}}{\partial\alpha_j}R_{c d}
+
R_{a b}\frac{\partial R_{c d}}{\partial\alpha_j}.
\]
This product-rule structure is the key point of the orientation derivative: once \(R\), \(\partial R/\partial\alpha_j\), \(T\), and \(\partial T/\partial\alpha_j\) are known, the sensitivity follows directly from the quadrature expression above.
